\documentclass[11pt, a4paper, twoside, titlepage]{book}

% Paquetes usados

\usepackage[galician]{babel}
\usepackage[a4paper, top=4cm, bottom=4cm, left=3.3cm, right=3.2cm]{geometry}
\usepackage[utf8]{inputenc}
\usepackage{graphicx}
\usepackage{fancyhdr}
\usepackage{subfigure}
\usepackage{booktabs}
\usepackage{url}
\usepackage[table,xcdraw]{xcolor}
\usepackage[official]{eurosym}
\usepackage{amsmath}

% Aumentar la separación entre líneas
\renewcommand{\baselinestretch}{1.5}
\graphicspath{ {figures/} }


\begin{document}


\pagestyle{empty}

\begin{titlepage}
\begin{center}

\includegraphics[scale=0.16]{logo.eps}

\vspace{0.5cm}
FACULTADE DE INFORMÁTICA

\vspace*{1cm}

\Large{TRABALLO FIN DE GRAO}

\Large{GRAO EN ENXEÑARÍA INFORMÁTICA}

Mención en Enxeñaría do Software

\vspace*{2cm}

\textbf{\LARGE{Aplicación web para a análise do comportamento humano en zoas transitadas}}

\end{center}

\vspace*{4cm}

\begin{flushright}
\large{
\textbf{Autor:} Xoan Iago Suárez Canosa\\
\textbf{Director:} Cancela Barizo, Brais \\
\textbf{Director:} González Penedo, Manuel Francisco \\
\textbf{Director:} Novo Buján, Jorge \\
\vspace{0.5cm}
A Coruña, \today}
\end{flushright}

\begin{flushleft}
 \includegraphics[scale=0.3]{logoOKI.png}
\end{flushleft}


\end{titlepage}

\cleardoublepage

% Numeración romana de páginas
\pagenumbering{roman}

% Definimos el encabezado y pie de página
\pagestyle{fancy}
\renewcommand{\chaptermark}[1]{%
\markboth{\thechapter.\ #1}{}}
\renewcommand{\sectionmark}[1]{%
\markright{\thesection.\ #1}{}}
\fancyhead{}
\fancyhead[LE,RO]{\thepage}
\fancyhead[LO,ER]{\rightmark}
\fancyfoot{}

\chapter{Introdución}
    \section{Descrición xeral da Aplicación}

        Este proxecto consiste nunha aplicación web que emprega funcionalidades relativas á análise do
        comportamento en zoas transitadas, incluíndo visualización de vídeo e distintas capas que
        mostran información de alto nivel sobre o comportamento detectado.\\

        En concreto esta ferramenta encargase de detectar a todos aqueles obxectos ou persoas que 
        aparecen nunha secuencia de vídeo, aplicando un algoritmo para medir o ``grao de anormalidade''
        da súa conduta en base aos movementos que realizan.\\

    \section{Obxectivos e vantaxes}
        Esta aplicación proporciona aos investigadores a vantaxe de poder probar os seus algoritmos de
        análise sobre un vídeo determinado, vendo os resultados a través dunha interface web e 
        quedando estes resultados almacenados e publicados para que calquera outra persoa que
        acceda á web poida visualizalos.
        
        Ao mesmo tempo tamén ofrece a posibilidade de que un usuario calquera queira probar as 
        capacidades do determinado algoritmo co que a aplicación estea a funcionar, podendo neste
        caso observar en que situacións o algoritmo funciona correctamente e en cales funciona de 
        forma errática.
        
    \section{Usuarios e Developers da Aplicación}
        O Developer principal da aplicación foi Iago Suárez, empregando como mostra un dos algoritmos
        desenvoltos polo laboratorio VARPA, mais dado que este é un proxecto de software libre 
        publicado, calquera outra persoa que o desexe pode descargar o software e facer sobre el as
        modificacións que considere. De feito esta aplicación está especialmente pensada para que 
        calquera investigador que estea a elaborar un algoritmo de análise de vídeo poida descargala
        e valerse de ela para amosar os resultados do seu traballo.
        
        Por outra parte os usuarios son todas aquelas persoas interesadas en coñecer as capacidades
        do algoritmo de análise que se estea a empregar.
        
\chapter{Requisitos de Hardware e Software}

    \section{Requisitos de Hardware}
        Non se especifican uns requisitos básicos, pero para o entorno de desenvolvemento é 
        recomendable polo menos dispor de unha máquina de 1GB de RAM, de un espazo libre en disco de 
        3 GB que poderá almacenar como mínimo 30 vídeos, procesador de alomenos 1.66 GHz e 
        obviamente conexión a internet para a descarga das dependencias precisas.

    \section{Requisitos de Software}
        En canto ao software é preciso un sistema operativo baseado en Linux, sendo recomendada por ser
        o entorno de desenvolvemento a versión 14.04 de Ubuntu. A linguaxe na que o servidor da 
        aplicación corre é Python 3.4, polo que compre tela instalada e asignada como versión de
        python por defecto.
        
\chapter{Procedementos de Instalación, arranque e finalización da Aplicación}

    \section{Contido do disco de instalación}
        \begin{itemize}
         \item \textbf{docs:} Documentación do proxecto
         \item \textbf{src:} Código fonte da aplicación
         \item \textbf{...:} Memoria da aplicación
         \item \textbf{...:} Manual de Usuario da aplicación
        \end{itemize}

    
    \section{Procedemento de Instalación}
        Todo o proceso de instalación está debidamente explicado nos ficheiros README.md
        e production-README.md, contendo no primeiro os pasos para a instalación no entorno
        de desenvolvemento e no segundo os pasos para o entorno de produción, ambos sobre unha
        máquina Ubuntu 14.04.
        
    \section{Arranque da Aplicación}
        Para arrancar o servidor da aplicación basta con:
        
        \begin{verbatim}
    cd src
    python manage.py migrate
    python manage.py createsuperuser --username=admin --email=admin@udc.es
    python manage.py runserver
        \end{verbatim}

        Co servidor arrincado xa se pode conectar ao porto local no que se iniciase ( normalmente 
        a dirección será http://127.0.0.1:8000/) como se ve na figura \ref{fig:Arranque}.
        
        \begin{figure}[htp]
        \begin{center}
            \includegraphics[scale=0.45]{figures/Arranque.png}
            \caption{Captura do arranque do servidor de desenvolvemento de Django.}
        \label{fig:Arranque}
        \end{center}
        \end{figure}
        
    \section{Finalización da Aplicación}
        Para deter o servidor chega con pulsar as teclas CONTROL-C sobre a terminal na que o estamos
        a executar. 

\chapter{Descrición de Menús, Pantallas e Procesos}
    \section{Menús, pantallas e procesos de la Aplicación}
        A aplicación está baseada nun menú central que se atopa na parte superior da aplicación como
        se pode ver na figura \ref{fig:menu} e dende o cal se pode acceder a tódalas partes da
        aplicación.
        
    \begin{figure}[htp]
    \begin{center}
        \includegraphics[scale=0.6]{figures/menu.png}
        \caption{Captura do arranque do servidor de desenvolvemento de Django.}
    \label{fig:menu}
    \end{center}
    \end{figure}
        
    \subsection{Páxina principal}
    
    Na páxina principal (figura \ref{fig:paxPrincipal}) vemos unha imaxe de benvida xunto cun breve 
    resumo do obxectivo da aplicación e de seguido unha lista dos último vídeos subidos.
    
        \begin{figure}[htp]
        \begin{center}
            \includegraphics[scale=0.4]{figures/paxPrincipal.png}
            \caption{Páxina principal da aplicación.}
        \label{fig:paxPrincipal}
        \end{center}
        \end{figure}

    \subsection{Listado de Vídeos}

    Se facemos clic no botón do menú ``Videos'' a aplicación conduciramos á páxina que contén o listado de 
    tódolos vídeos da aplicación da imaxe \ref{fig:paxVideos}. Nela aparecen unha lista de vídeos 
    ordenados por orde de subida e paxinados de cinco en cinco. Na mesma páxina podemos realizar
    buscas segundo o nome do vídeo.
    
        \begin{figure}[htp]
        \begin{center}
            \includegraphics[scale=0.6]{figures/paxVideos.png}
            \caption{Páxina de Vídeos da aplicación.}
        \label{fig:paxVideos}
        \end{center}
        \end{figure}

    \subsection{Reprodución de vídeos e análise}
    
    Se facemos clic sobre algún vídeo da aplicación (tanto da lista de vídeos como da páxina 
    principal) o seu enlace levaranos a esta páxina na que podemos reproducilo. A maiores 
    superpóñense sobre o vídeo unha serie de capas de información que podemos configurar co panel de
    control ``Controls'' como na figura \ref{fig:paxDetails}, nel veremos os seguintes botóns:
    
    \begin{itemize}
     \item \textbf{Show Analysis:} Amosa ou oculta todas as capas de información sobre o vídeo.
     \item \textbf{Use different colors:} Habilitase ou deshabitase a identificación de obxectos
        mediante diferentes cores.
     \item \textbf{Show all detections:} Amosa ou oculta unha táboa na que se ven tódalas deteccións
        atopadas no vídeo.
     \item \textbf{Filter by abnormality rate:} Desprega un sub-panel que permite filtrar as 
        deteccións en función do anómalo que sexa o seu comportamento. Se unha detección é 
        considerada como sospeitosa marcarase en cor vermella, de ser de confianza en cor azul e por
        último se inda non foi analizada de cor negra. De estar activado o subpanel
        poderemos ver:
        \begin{itemize}
         \item \textbf{Barra Selectora:} Nela podemos seleccionar a partir de que valor 
            consideraremos anómala unha traxectoria.
         \item \textbf{Entrada de Texto:} Nela ao igual que na barra selectora podemos seleccionar 
            a partir de que valor consideraremos anómala unha traxectoria.
         \item \textbf{Throw Pop-ups:} Este botón empregase para indicar que se lancen ou non ventás
            emerxentes ao detectar un obxecto sospeitoso.
        \end{itemize}
    \end{itemize}

        \begin{figure}[htp]
        \begin{center}
            \includegraphics[scale=0.4]{figures/paxDetails.png}
            \caption{Páxina de reprodución de vídeos.}
        \label{fig:paxDetails}
        \end{center}
        \end{figure}
        
    \subsection{Subida de vídeos}
    \label{sec:subidaVideo}
        Unha vez logueados, se facemos clic sobre o nome do noso usuario na marxe superior dereita 
        despregase unha lista e opcións onde a primeira de elas será ``Upload new video'', facendo
        clic nela accederemos a esta páxina para a subida de vídeos (imaxe \ref{fig:paxUpload}). No
        formulario que aparece a continuación poderemos subir un vídeo de menos de 30 MB nun dos
        seguintes formatos:
        \begin{itemize}
         \item .mp4
         \item .ogv
         \item .avi
         \item .flv
         \item .mpg 
        \end{itemize}

        \begin{figure}[htp]
        \begin{center}
            \includegraphics[scale=0.5]{figures/paxUpload.png}
            \caption{Páxina de subida de vídeos.}
        \label{fig:paxUpload}
        \end{center}
        \end{figure} 
        
    \subsection{Páxina sobre a aplicación}
        Se facemos clic sobre o botón do menú principal que pon ``About'' poderemos acceder a unha 
        páxina que nos informa sobre a aplicación e o xeito no que está construída como se ve na 
        figura \ref{fig:paxAbout}.
        
        \begin{figure}[htp]
        \begin{center}
            \includegraphics[scale=0.4]{figures/paxAbout.png}
            \caption{Páxina sobre a aplicación (About).}
        \label{fig:paxAbout}
        \end{center}
        \end{figure}
        
    \subsection{Páxina de Administración}
        
        A paxina de administración serve para que o administrados do sitio web poida borrar, editar
        ou crear novas entidades da aplicación como poden ser por exemplo usuarios ou vídeos. A 
        páxina luce como se ve na imaxe \ref{fig:paxAdmin}.
        
        \begin{figure}[htp]
        \begin{center}
            \includegraphics[scale=0.55]{figures/paxAdmin.png}
            \caption{Páxina de administración da aplicación.}
        \label{fig:paxAdmin}
        \end{center}
        \end{figure}
        
\chapter{Descrición de Procedementos do Usuario}

    \section{Crear un usuario}
    Para crear un usuario chega con cubrir o formulario \ref{fig:register} que aparece na páxina principal.
    
    \begin{figure}[htp]
    \begin{center}
        \includegraphics[scale=0.6]{figures/register.png}
        \caption{Páxina de administración da aplicación.}
    \label{fig:register}
    \end{center}
    \end{figure}

    \section{Subir un vídeo}
    
    Para subir un vídeo primeiro accederemos á paxina de subida como se explicou no apartado 
    \ref{sec:subidaVideo}, unha vez aí cubrimos o vídeo e seleccionamos ``Submit'', este vídeo comezará
    a subirse informándonos do seu proceso mediante unha notificación. Cando a todos os procesos de
    subida rematen (Carga, conversión, análise e extracción de imaxes) apareceranos unha 
    notificación cun botón como o que se ve na figura \ref{fig:notFinished}.
    
    \begin{figure}[htp]
    \begin{center}
        \includegraphics[scale=0.6]{figures/notFinished.png}
        \caption{Notificación de que a subida do vídeo rematou.}
    \label{fig:notFinished}
    \end{center}
    \end{figure}
    
    Facendo clic nese botón de ``Finish'' a aplicación conduciranos á unha páxina na que se 
    seleccionará a imaxe representativa do vídeo como se amosa na imaxe 
    \ref{fig:SuccessfulUploadScreen}, cando teñamos seleccionada esta imaxe e premamos no boton de 
    ``Submit'' a aplicación levaranos á páxina para a visualización de vídeo na que poderemos ver o 
    resultado da análise.
    
    \begin{figure}[htp]
    \begin{center}
        \includegraphics[scale=0.35]{figures/SuccessfulUploadScreen.png}
        \caption{Captura de pantalla da páxina web SuccessfulUpload.}
    \label{fig:SuccessfulUploadScreen}
    \end{center}
    \end{figure}
    
    
    \section{Visualizar un vídeo e a súa análise}
    
    Para reproducir un vídeo o proceso a seguir é primeiro buscar o vídeo que se quere visualizar, 
    facer clic sobre el para acceder á páxina de reprodución de vídeo e unha vez aí pulsar sobre o
    elemento vídeo ou sobre a barra espaciadora para que o vídeo comeze a reproducirse.
    
    
% Numeración normal de páginas
\pagenumbering{arabic}


% Definimos el encabezado y pie de página
\fancyhead{}
\fancyhead[LE,RO]{\thepage}
\fancyhead[LO]{\rightmark}
\fancyhead[ER]{\leftmark}

% Definimos el encabezado y pie de página
\fancyhead{}
\fancyhead[LE,RO]{\thepage}
\fancyhead[LO,ER]{\rightmark}

% Glosario de términos
%\printglossary

\end{document}
